% Beamer slides on Dephasing Time T2 for Graduate Physics Students
\documentclass{beamer}
\usetheme{Madrid}
\usepackage{amsmath, amssymb, physics, graphicx}
\title{Dephasing Time $T_2$ in Quantum Systems}
\author{FYS5419}
\date{}

\begin{document}

\begin{frame}\titlepage\end{frame}

  % Slide 1
  \begin{frame}{What is Dephasing?}
    Dephasing refers to the loss of phase coherence between quantum states due to interactions with the environment.
    \begin{itemize}
    \item Impacts off-diagonal elements of the density matrix.
    \item Leads to reduced visibility in interference experiments.
    \item A key limiting factor in quantum computing, sensing, and communication.
    \end{itemize}
  \end{frame}

  % Slide 2
  \begin{frame}{Density Matrix Perspective}
    For a qubit in state $\rho$:
    \rho = \begin{pmatrix} \rho_{00} & \rho_{01} \\ \rho_{10} & \rho_{11} \end{pmatrix}
    Dephasing primarily suppresses the coherence terms:
    \rho_{01}(t) = \rho_{01}(0)e^{-t/T_2}
  \end{frame}

  % Slide 3
  \begin{frame}{Bloch Sphere Representation}
    \begin{itemize}
    \item Dephasing corresponds to shrinking of the Bloch vector in the transverse plane.
    \item Longitudinal component (along $z$) is unaffected.
    \end{itemize}
    (Insert Bloch sphere illustration)
  \end{frame}

  % Slide 4
  \begin{frame}{$T_1$ vs $T_2$ Times}
    \begin{itemize}
    \item $T_1$: Energy relaxation (population decay)
    \item $T_2$: Phase relaxation (coherence decay)
    \item Always: $\frac{1}{T_2} = \frac{1}{2T_1} + \frac{1}{T_\phi}$
    \end{itemize}
  \end{frame}

  % Slide 5
  \begin{frame}{Sources of Dephasing}
    \begin{itemize}
    \item Magnetic field fluctuations
    \item Charge noise
    \item Spin baths
    \item Phonons and lattice vibrations
    \item Classical low-frequency noise ($1/f$ noise)
    \end{itemize}
  \end{frame}

  % Slide 6
  \begin{frame}{Noise Models}
    Common noise spectral densities:
    \begin{itemize}
    \item White noise: constant spectral density
    \item $1/f$ noise: dominates at low frequencies
    \item Ohmic noise: linear in frequency
    \item Lorentzian noise: due to two-level fluctuators
    \end{itemize}
  \end{frame}

  % Slide 7
  \begin{frame}{Master Equation Formulation}
    Lindblad equation with dephasing:
    \dot{\rho} = -i[H,\rho] + \gamma_\phi (\sigma_z \rho \sigma_z - \rho)
    Here:
    T_2 = \frac{1}{2\gamma_\phi}
  \end{frame}

  % Slide 8
  \begin{frame}{Ramsey Experiment}
    Ramsey sequence measures free-induction decay:
    \begin{itemize}
    \item Prepare superposition
    \item Free evolution time $t$
    \item Final $\pi/2$ pulse
    \end{itemize}
    Signal:
    S(t) = S_0 e^{-t/T_2} \cos(\Delta t)
  \end{frame}

  % Slide 9
  \begin{frame}{Spin Echo}
    Spin echo recovers coherence lost due to slow noise.
    S(t) = S_0 e^{-t/T_2^{\rm echo}}
    Typically: $T_2^{\rm echo} > T_2$.
  \end{frame}

  % Slide 10
  \begin{frame}{$T_2^$ Time}
    \begin{itemize}
    \item $T_2^$ includes inhomogeneous broadening.
    \item Extracted from Ramsey without echo.
    \item Relation: $T_2^* < T_2 < T_2^{\rm echo}$.
    \end{itemize}
  \end{frame}

  % Slide 11
  \begin{frame}{Typical $T_2$ Values: Superconducting Qubits}
    Examples:
    \begin{itemize}
    \item Transmons: $T_2 \sim 50$–$200~\mu$s
    \item Flux qubits: shorter due to flux noise
    \end{itemize}
    (Insert chart)
  \end{frame}

  % Slide 12
  \begin{frame}{Typical $T_2$ Values: Trapped Ions}
    \begin{itemize}
    \item $T_2$ times often exceed seconds
    \item Limited mainly by magnetic field drift
    \end{itemize}
  \end{frame}

  % Slide 13
  \begin{frame}{Typical $T_2$ Values: NV Centers in Diamond}
    \begin{itemize}
    \item Room-temperature $T_2$: $100$–$500~\mu$s
    \item With isotopic purification: up to ms
    \end{itemize}
  \end{frame}

  % Slide 14
  \begin{frame}{Typical $T_2$ Values: Semiconductor Spin Qubits}
    \begin{itemize}
    \item $T_2^*$ often ns to µs range
    \item Echo-improved $T_2$: tens to hundreds of µs
    \end{itemize}
  \end{frame}

  % Slide 15
  \begin{frame}{Mathematical Derivation: Free Induction Decay}
    Given Hamiltonian with fluctuating frequency:
    H = \frac{1}{2}(\omega_0 + \delta \omega(t))\sigma_z
    Phase:
    \phi(t) = \int_0^t dt’ \, \delta\omega(t’)
    Coherence:
    \langle e^{i\phi(t)} \rangle = e^{-t/T_2}
  \end{frame}

  % Slide 16
  \begin{frame}{Noise Correlation Functions}
    If noise is Gaussian:
    T_2^{-1} = \int_0^\infty d\omega \, S(\omega)
    where $S(\omega)$ is the noise spectrum.
  \end{frame}

  % Slide 17
  \begin{frame}{Experimental Signatures of Dephasing}
    \begin{itemize}
    \item Linewidth broadening
    \item Amplitude decay in Ramsey/echo experiments
    \item Reduced gate fidelity
    \end{itemize}
  \end{frame}

  % Slide 18
  \begin{frame}{Role of Dephasing in Quantum Computing}
    \begin{itemize}
    \item Limits coherence time window for gates
    \item Fundamental constraint on error correction
    \item Determines clock speed of logical operations
    \end{itemize}
  \end{frame}

  % Slide 19
  \begin{frame}{Mitigating Dephasing}
    \begin{itemize}
    \item Dynamical decoupling sequences
    \item Material purification
    \item Cooling and shielding
    \item Optimal control pulses
    \item Noise spectroscopy and adaptive correction
    \end{itemize}
  \end{frame}

  % Slide 20
  \begin{frame}{Summary}
    \begin{itemize}
    \item $T_2$ characterizes phase coherence
    \item Essential for quantum technologies
    \item Strong platform dependence
    \item Experimental control and theory needed for mitigation
    \end{itemize}
  \end{frame}

  % –– Additional slides to extend duration to ~2 hours ––

  % Slide: Advanced derivation of chi(t)
  \begin{frame}{Advanced Derivation of $\chi(t)$}
    Starting from the cumulant expansion for Gaussian noise:
    \chi(t)=\frac{1}{2}\int_0^t dt_1\int_0^t dt_2 \, C(t_1-t_2),
    where $C(\tau)$ is the noise autocorrelation. Using its spectral density:
    C(\tau)=\frac{1}{2\pi}\int_{-\infty}^{\infty} d\omega \, S(\omega)e^{-i\omega\tau}.
    Combining yields the standard filter function representation.
  \end{frame}

  % Slide: Filter function formalism
  \begin{frame}{Filter Function Formalism}
    A control sequence defines a modulation function $y(t)$, leading to:
    \chi(t)=\frac{1}{2\pi}\int_0^{\infty} d\omega \, S(\omega) \frac{|F(\omega)|^2}{\omega^2},
    where $F(\omega)=\int_0^{t} dt’ , y(t’) e^{i\omega t’}$.
    \begin{itemize}
    \item Ramsey: $y(t)=1$.
    \item Spin echo: $y(t)=1$ for $t<\tau$, $-1$ afterwards.
    \end{itemize}
  \end{frame}

  % Slide: Comparison of control sequences
  \begin{frame}{Comparison of Control Sequences}
    \begin{itemize}
    \item Ramsey: sensitive to low-frequency noise.
    \item Hahn Echo: cancels quasi-static noise.
    \item CPMG: multiple $\pi$ pulses increase suppression bandwidth.
    \end{itemize}
    Include conceptual plots of $|F(\omega)|^2$.
  \end{frame}

  % Slide: Dephasing in superconducting qubits (materials)
  \begin{frame}{Materials Origins of Dephasing in Superconducting Qubits}
    Sources:
    \begin{itemize}
    \item Two-level fluctuators in amorphous oxides.
    \item Flux noise from surface spins.
    \item Charge noise from dielectric loss.
    \end{itemize}
    Mitigations include surface cleaning, substrate engineering, and improved fabrication.
  \end{frame}

  % Slide: Dephasing in trapped ions
  \begin{frame}{Dephasing in Trapped-Ion Systems}
    Dominant mechanisms:
    \begin{itemize}
    \item Magnetic field drifts.
    \item Laser phase fluctuations.
    \item Motional heating.
    \end{itemize}
    Use of sympathetic cooling and clock transitions minimizes noise.
  \end{frame}

  % Slide: Dephasing in NV centers
  \begin{frame}{NV Centers: Role of Spin Bath}
    The $^{13}$C nuclear spin bath causes spectral diffusion.
    \begin{itemize}
    \item Echo improves $T_2$ by refocusing bath dynamics.
    \item Isotopic purification reduces nuclear spin density.
    \end{itemize}
  \end{frame}

  % Slide: Semiconductor spins detailed
  \begin{frame}{Semiconductor Spin Qubits: Charge Noise and Hyperfine Effects}
    Key mechanisms:
    \begin{itemize}
    \item Hyperfine fluctuations (GaAs).
    \item Charge noise modifying exchange energy.
    \end{itemize}
    Silicon qubits mitigate noise via isotopic purification and optimized gate stacks.
  \end{frame}

  % Slide: Full table of T1/T2 across systems
  \begin{frame}{Extended Table: $T_1$, $T_2^*$, and $T_2$ Across Platforms}
    Include table with:
    \begin{itemize}
    \item Superconducting qubits.
    \item Spin qubits.
    \item NV centers.
    \item Trapped ions.
    \item Neutral atoms.
    \item Donor spins.
    \end{itemize}
  \end{frame}

  % Slide: Experimental extraction details
  \begin{frame}{Extracting Coherence Times: Practical Considerations}
    \begin{itemize}
    \item Fit models: exponential, stretched exponential, Gaussian.
    \item Frequency detuning corrections.
    \item Accounting for SPAM errors.
    \end{itemize}
  \end{frame}

  % Slide: Noise spectroscopy derivation
  \begin{frame}{Derivation: Noise Spectroscopy from CPMG}
    Peak sensitivity at
    \omega_n \approx \frac{\pi n}{t},
    allowing sampling of $S(\omega)$ across frequencies.
  \end{frame}

  % Slide: Bloch-Redfield theory
  \begin{frame}{Bloch–Redfield Theory}
    Redfield tensor gives dephasing and relaxation rates:
    \Gamma_{ij,kl}=\sum_{\alpha\beta} R^{\alpha\beta}{ij,kl} S{\alpha\beta}(\omega).
    Used to model multi-level systems.
  \end{frame}

  % Slide: Quantum error correction relevance
  \begin{frame}{Relevance for Quantum Error Correction}
    \begin{itemize}
    \item Coherence time sets maximum logical gate depth.
    \item Fault-tolerant thresholds require low dephasing rates.
    \item Strategies depend on dominant noise channel.
    \end{itemize}
  \end{frame}

  % Slide: Examples of $T_2$ improvements
  \begin{frame}{Documented Improvements in $T_2$}
    \begin{itemize}
    \item 3D transmons: quality factors improved by orders of magnitude.
    \item Donor spins: $T_2$ up to 100 s.
    \item NV centers: millisecond-scale coherence.
    \end{itemize}
  \end{frame}

  % Slide: Experimental Bloch sphere animations (conceptual)
  \begin{frame}{Visualizing Dephasing on the Bloch Sphere}
    Show successive shrinking of transverse radius.
    Discuss connection to noise bandwidth.
  \end{frame}

  % Slide: Numerical simulations
  \begin{frame}{Simulating Dephasing}
    \begin{itemize}
    \item Classical noise generation (Ornstein–Uhlenbeck, $1/f$ noise).
    \item Solving stochastic Schrödinger equation.
    \item Averaging over many trajectories.
    \end{itemize}
  \end{frame}

  % Slide: Open quantum systems overview
  \begin{frame}{Open Quantum Systems: Overview}
    \begin{itemize}
    \item Markovian vs non-Markovian dynamics.
    \item Master equations.
    \item Quantum trajectories.
    \end{itemize}
  \end{frame}

  % Slide: Experimental platforms summary
  \begin{frame}{Platform Comparison Summary}
    Evaluate:
    \begin{itemize}
    \item Dominant noise sources.
    \item Control techniques.
    \item Typical coherence times.
    \item Scalability considerations.
    \end{itemize}
  \end{frame}

  % –– Additional slides to extend duration to ~2 hours ––

  % Slide: Advanced derivation of chi(t)
  \begin{frame}{Advanced Derivation of $\chi(t)$}
    Starting from the cumulant expansion for Gaussian noise:
    \chi(t)=\frac{1}{2}\int_0^t dt_1\int_0^t dt_2 \, C(t_1-t_2),
    where $C(\tau)$ is the noise autocorrelation. Using its spectral density:
    C(\tau)=\frac{1}{2\pi}\int_{-\infty}^{\infty} d\omega \, S(\omega)e^{-i\omega\tau}.
    Combining yields the standard filter function representation.
  \end{frame}

  % Slide: Filter function formalism
  \begin{frame}{Filter Function Formalism}
    A control sequence defines a modulation function $y(t)$, leading to:
    \chi(t)=\frac{1}{2\pi}\int_0^{\infty} d\omega \, S(\omega) \frac{|F(\omega)|^2}{\omega^2},
    where $F(\omega)=\int_0^{t} dt’ , y(t’) e^{i\omega t’}$.
    \begin{itemize}
    \item Ramsey: $y(t)=1$.
    \item Spin echo: $y(t)=1$ for $t<\tau$, $-1$ afterwards.
    \end{itemize}
  \end{frame}

  % Slide: Comparison of control sequences
  \begin{frame}{Comparison of Control Sequences}
    \begin{itemize}
    \item Ramsey: sensitive to low-frequency noise.
    \item Hahn Echo: cancels quasi-static noise.
    \item CPMG: multiple $\pi$ pulses increase suppression bandwidth.
    \end{itemize}
    Include conceptual plots of $|F(\omega)|^2$.
  \end{frame}

  % Slide: Dephasing in superconducting qubits (materials)
  \begin{frame}{Materials Origins of Dephasing in Superconducting Qubits}
    Sources:
    \begin{itemize}
    \item Two-level fluctuators in amorphous oxides.
    \item Flux noise from surface spins.
    \item Charge noise from dielectric loss.
    \end{itemize}
    Mitigations include surface cleaning, substrate engineering, and improved fabrication.
  \end{frame}

  % Slide: Dephasing in trapped ions
  \begin{frame}{Dephasing in Trapped-Ion Systems}
    Dominant mechanisms:
    \begin{itemize}
    \item Magnetic field drifts.
    \item Laser phase fluctuations.
    \item Motional heating.
    \end{itemize}
    Use of sympathetic cooling and clock transitions minimizes noise.
  \end{frame}

  % Slide: Dephasing in NV centers
  \begin{frame}{NV Centers: Role of Spin Bath}
    The $^{13}$C nuclear spin bath causes spectral diffusion.
    \begin{itemize}
    \item Echo improves $T_2$ by refocusing bath dynamics.
    \item Isotopic purification reduces nuclear spin density.
    \end{itemize}
  \end{frame}

  % Slide: Semiconductor spins detailed
  \begin{frame}{Semiconductor Spin Qubits: Charge Noise and Hyperfine Effects}
    Key mechanisms:
    \begin{itemize}
    \item Hyperfine fluctuations (GaAs).
    \item Charge noise modifying exchange energy.
    \end{itemize}
    Silicon qubits mitigate noise via isotopic purification and optimized gate stacks.
  \end{frame}

  % Slide: Full table of T1/T2 across systems
  \begin{frame}{Extended Table: $T_1$, $T_2^*$, and $T_2$ Across Platforms}
    Include table with:
    \begin{itemize}
    \item Superconducting qubits.
    \item Spin qubits.
    \item NV centers.
    \item Trapped ions.
    \item Neutral atoms.
    \item Donor spins.
    \end{itemize}
  \end{frame}

  % Slide: Experimental extraction details
  \begin{frame}{Extracting Coherence Times: Practical Considerations}
    \begin{itemize}
    \item Fit models: exponential, stretched exponential, Gaussian.
    \item Frequency detuning corrections.
    \item Accounting for SPAM errors.
    \end{itemize}
  \end{frame}

  % Slide: Noise spectroscopy derivation
  \begin{frame}{Derivation: Noise Spectroscopy from CPMG}
    Peak sensitivity at
    \omega_n \approx \frac{\pi n}{t},
    allowing sampling of $S(\omega)$ across frequencies.
  \end{frame}

  % Slide: Bloch-Redfield theory
  \begin{frame}{Bloch–Redfield Theory}
    Redfield tensor gives dephasing and relaxation rates:
    \Gamma_{ij,kl}=\sum_{\alpha\beta} R^{\alpha\beta}{ij,kl} S{\alpha\beta}(\omega).
    Used to model multi-level systems.
  \end{frame}

  % Slide: Quantum error correction relevance
  \begin{frame}{Relevance for Quantum Error Correction}
    \begin{itemize}
    \item Coherence time sets maximum logical gate depth.
    \item Fault-tolerant thresholds require low dephasing rates.
    \item Strategies depend on dominant noise channel.
    \end{itemize}
  \end{frame}

  % Slide: Examples of $T_2$ improvements
  \begin{frame}{Documented Improvements in $T_2$}
    \begin{itemize}
    \item 3D transmons: quality factors improved by orders of magnitude.
    \item Donor spins: $T_2$ up to 100 s.
    \item NV centers: millisecond-scale coherence.
    \end{itemize}
  \end{frame}

  % Slide: Experimental Bloch sphere animations (conceptual)
  \begin{frame}{Visualizing Dephasing on the Bloch Sphere}
    Show successive shrinking of transverse radius.
    Discuss connection to noise bandwidth.
  \end{frame}

  % Slide: Numerical simulations
  \begin{frame}{Simulating Dephasing}
    \begin{itemize}
    \item Classical noise generation (Ornstein–Uhlenbeck, $1/f$ noise).
    \item Solving stochastic Schrödinger equation.
    \item Averaging over many trajectories.
    \end{itemize}
  \end{frame}

  % Slide: Open quantum systems overview
  \begin{frame}{Open Quantum Systems: Overview}
    \begin{itemize}
    \item Markovian vs non-Markovian dynamics.
    \item Master equations.
    \item Quantum trajectories.
    \end{itemize}
  \end{frame}

  % Slide: Experimental platforms summary
  \begin{frame}{Platform Comparison Summary}
    Evaluate:
    \begin{itemize}
    \item Dominant noise sources.
    \item Control techniques.
    \item Typical coherence times.
    \item Scalability considerations.
    \end{itemize}
  \end{frame}

  % –– Additional slides to extend duration to ~2 hours ––

  % Slide: Advanced derivation of chi(t)
  \begin{frame}{Advanced Derivation of $\chi(t)$}
    Starting from the cumulant expansion for Gaussian noise:
    \[
    \\chi(t)=\\frac{1}{2}\\int_0^t dt_1\\int_0^t dt_2 \, C(t_1-t_2),
    \\]
      where $C(\tau)$ is the noise autocorrelation. Using its spectral density:
      \[
      C(\\tau)=\\frac{1}{2\\pi}\\int_{-\\infty}^{\\infty} d\\omega \, S(\\omega)e^{-i\\omega\\tau}.
      \\]
        Combining yields the standard filter function representation.
  \end{frame}

  % Slide: Filter function formalism
  \begin{frame}{Filter Function Formalism}
    A control sequence defines a modulation function $y(t)$, leading to:
    \[
    \\chi(t)=\\frac{1}{2\\pi}\\int_0^{\\infty} d\\omega \, S(\\omega) \\frac{|F(\\omega)|^2}{\\omega^2},
    \\]
      where $F(\omega)=\int_0^{t} dt’ , y(t’) e^{i\omega t’}$.
      \begin{itemize}
      \item Ramsey: $y(t)=1$.
      \item Spin echo: $y(t)=1$ for $t<\tau$, $-1$ afterwards.
      \end{itemize}
  \end{frame}

  % Slide: Comparison of control sequences
  \begin{frame}{Comparison of Control Sequences}
    \begin{itemize}
    \item Ramsey: sensitive to low-frequency noise.
    \item Hahn Echo: cancels quasi-static noise.
    \item CPMG: multiple $\pi$ pulses increase suppression bandwidth.
    \end{itemize}
  \end{frame}

  % Slide: Dephasing in superconducting qubits
  \begin{frame}{Materials Origins of Dephasing in Superconducting Qubits}
    Sources:
    \begin{itemize}
    \item Two-level fluctuators in amorphous oxides.
    \item Flux noise from surface spins.
    \item Charge noise from dielectric loss.
    \end{itemize}
    Mitigations include surface cleaning, substrate engineering, and improved fabrication.
  \end{frame}

  % Slide: Dephasing in trapped ions
  \begin{frame}{Dephasing in Trapped-Ion Systems}
    Dominant mechanisms:
    \begin{itemize}
    \item Magnetic field drifts.
    \item Laser phase fluctuations.
    \item Motional heating.
    \end{itemize}
    Use of sympathetic cooling and clock transitions minimizes noise.
  \end{frame}

  % Slide: Dephasing in NV centers
  \begin{frame}{NV Centers: Role of Spin Bath}
    The $^{13}$C nuclear spin bath causes spectral diffusion.
    \begin{itemize}
    \item Echo improves $T_2$ by refocusing bath dynamics.
    \item Isotopic purification reduces nuclear spin density.
    \end{itemize}
  \end{frame}

  % Slide: Semiconductor spins
  \begin{frame}{Semiconductor Spin Qubits: Charge Noise and Hyperfine Effects}
    Key mechanisms:
    \begin{itemize}
    \item Hyperfine fluctuations (GaAs).
    \item Charge noise modifying exchange energy.
    \end{itemize}
    Silicon qubits mitigate noise via isotopic purification and optimized gate stacks.
  \end{frame}

  % Slide: Extended coherence table
  \begin{frame}{Extended Table: $T_1$, $T_2^*$, $T_2$ Across Platforms}
    Additional table can include detailed numerical ranges and citations.
  \end{frame}

  % Slide: Extraction details
  \begin{frame}{Extracting Coherence Times: Practical Considerations}
    \begin{itemize}
    \item Fit models: exponential, stretched exponential, Gaussian.
    \item Frequency detuning corrections.
    \item Accounting for SPAM errors.
    \end{itemize}
  \end{frame}

  % Slide: Noise spectroscopy
  \begin{frame}{Noise Spectroscopy from CPMG}
    Peak sensitivity at $\omega_n \approx \pi n/t$, enabling sampling of $S(\omega)$.
  \end{frame}

  % Slide: Bloch-Redfield
  \begin{frame}{Bloch–Redfield Theory}
    Redfield tensor gives dephasing and relaxation rates. Applicable to multi-level systems.
  \end{frame}

  % Slide: QEC relevance
  \begin{frame}{Relevance for Quantum Error Correction}
    Coherence limits gate depth and error thresholds; noise characterization guides code design.
  \end{frame}

  % Slide: Improvements
  \begin{frame}{Documented Improvements in $T_2$}
    Examples from transmons, donor spins, NV centers.
  \end{frame}

  % Slide: Visualization
  \begin{frame}{Visualizing Dephasing on Bloch Sphere}
    Successive shrinking of transverse radius shows decoherence progression.
  \end{frame}

  % Slide: Simulation
  \begin{frame}{Simulating Dephasing}
    Classical noise generation, stochastic Schrödinger integration, ensemble averaging.
  \end{frame}

  % Slide: Open quantum systems
  \begin{frame}{Open Quantum Systems Overview}
    Markovian vs non-Markovian dynamics; master equations; quantum trajectories.
  \end{frame}

  % Slide: Platform summary
  \begin{frame}{Platform Comparison Summary}
    Evaluate noise sources, coherence, control, scalability.
  \end{frame}

  \end{document}
